\section{Repair under a Budget}
\label{sec:repair}

\subsection{Design}
\label{sec:repair-design}

If only some lost links can be restored, which should come first? We removed
30\% or 50\% of the edges at random, using ten damage orders, and allowed each
policy to restore 5\%, 10\% or 20\% of the deleted edges. At 50\% damage these
budgets are 2,947, 5,895 and 11,789 edges. Every policy starts from the same
damaged graphs and restores a prefix of its own ranking of the deleted edges.
\begin{itemize}[leftmargin=*,itemsep=2pt]
\item \textbf{Random}: uniform order; the baseline.
\item \textbf{Edge betweenness}: deleted edges in order of their betweenness in
the intact graph, estimated from 100 sampled sources. This is structural repair and uses no labels.
\item \textbf{Relation type}: deleted drug--protein and disease--protein edges
only, in random order. This is layer-aware repair: it restores the two relation
layers that a two-hop route needs.
\item \textbf{Static task-aware} (labeled task-aware in the figures): each deleted edge scores one point for every
broken two-hop route of a training query that it alone would complete. Edges
are ranked by score, with ties broken at random.
\end{itemize}
Gains are measured on 60 held-out queries, split from the 140 training queries
by query identifier, and reported as paired differences from the unrepaired
state of the same damage order.

Two limitations of this design shape the results. First, the task-aware score
is static and saturates almost immediately. Only 31--46 deleted edges have a
positive score at 30\% damage, and 62--87 at 50\%, which is less than 3\% of
even the smallest budget. The rest of its budget is therefore spent on random
ties, three quarters of them protein--protein edges. Second, training and test
queries share drugs, so the held-out set is not independent of the training
routes. The task-aware results therefore show what a small number of
well-chosen edges can do, not the limit of task-aware repair.

\subsection{What the theory predicts}
\label{sec:repair-theory}

Equation~\eqref{eq:r2} predicts two of the policies without fitting. Random
repair of a share $b$ of the deleted edges raises the edge retention from $p$
to $p+b(1-p)$. Relation-type repair leaves protein--protein edges unchanged, so
it does not affect $R_2$ through them. It raises the retention of the 27,704
drug--protein and disease--protein edges instead, and at a 20\% budget after
50\% damage it restores 85\% of those that were deleted. Substituting these
retentions into Eq.~\eqref{eq:r2} for the test queries gives the predictions in
Table~\ref{tab:repair}.

Eleven of the twelve predictions (two policies, two damage levels, three
budgets) are slightly above the measured gains, by at most 1.5 standard errors
of the ten-order mean, and all are within one standard deviation. The twelve
comparisons share the same damage orders, so their common sign is not
independent evidence. It has a simple cause. Their unrepaired test-set success is 42.3\% at
$f=0.5$, against 40.2\% expected from Eq.~\eqref{eq:r2}, so these orders left
fewer failed queries to repair (Section~\ref{sec:blind-uncertainty}).

\begin{table}[!htbp]
\centering
\small
\caption{Gains in percentage points after restoring 20\% of the edges deleted at
$f=0.5$ (mean $\pm$ SD over ten damage orders). Unrepaired values are 42.3\%
($R_2$), 83.5\% ($R_3$) and 98.2\% ($S$). The predicted column uses
Eq.~\eqref{eq:r2}, and $S$ gains are in percentage points of nodes.}
\label{tab:repair}
\begin{tabular}{@{}lcccc@{}}
\toprule
Policy & $R_2$ predicted & $R_2$ measured & $R_3$ measured & $S$ measured \\
\midrule
Random            & 12.3 & $13.2 \pm 2.8$ & $8.3 \pm 5.8$  & $0.58 \pm 0.11$ \\
Edge betweenness  & --   & $42.0 \pm 4.0$ & $16.5 \pm 7.1$ & $1.78 \pm 0.13$ \\
Relation type     & 51.6 & $48.5 \pm 6.7$ & $15.7 \pm 7.3$ & $0.61 \pm 0.13$ \\
Static task-aware & --   & $24.8 \pm 6.3$ & $12.3 \pm 6.6$ & $0.67 \pm 0.09$ \\
\bottomrule
\end{tabular}
\end{table}

\subsection{Which repair works for which task}
\label{sec:repair-results}

Figure~\ref{fig:repair} shows the gains for all budgets. We compare policies
with paired $t$-tests over the ten damage orders.
\begin{itemize}[leftmargin=*,itemsep=2pt]
\item \textbf{Two hops, large budget.} Relation-type repair is best. At a 20\%
budget it beats edge betweenness by 3.8 points at $f=0.3$ ($P=0.006$, 8 of 10
orders) and by 6.5 points at $f=0.5$ ($P=0.002$, 9 of 10). Only the second
survives a Bonferroni correction for the 24 comparisons we ran. Edge betweenness
spends 68\% of its budget on protein--protein edges, which cannot help a
two-hop route (Figure~\ref{fig:repair-types}).
\item \textbf{Two hops, small budget.} The advantage disappears. At a 5\%
budget, edge betweenness is ahead by 4.0 points at $f=0.3$ ($P=0.04$, not
significant after correction), and the two are indistinguishable at $f=0.5$. Drug links
are often bridges, so edge betweenness picks them early: they make up 1.3\% of
its first repairs, three times their share under random repair. The task-aware
score is also competitive here, because its few positive-score edges come
first.
\item \textbf{Three hops.} Edge betweenness is best or tied. At a 5\% budget it
beats relation-type repair by 4.0 points at $f=0.3$ ($P=0.003$) and by 11.0 points
at $f=0.5$ ($P=0.001$, which survives the correction). At a 20\% budget after 50\% damage it restores $R_3$ to 100\% in
every order, so the remaining differences are capped. Once one PPI step is allowed,
protein--protein bridges become useful.
\item \textbf{Connectivity.} The giant component has only 1.8 points of headroom
at $f=0.5$. Edge
betweenness recovers all of it by a 10\% budget. The other policies recover
18--25\% of it at a 10\% budget and 33--38\% at 20\%.
\end{itemize}

\begin{figure}[!htbp]
  \centering
  \includegraphics[width=\linewidth]{act3/final/primekg_5000_bfs_v1/fig_restoration.pdf}
  \caption{Held-out retrieval gains for the four policies at $L=2$ and $L=3$, after
  30\% or 50\% edge loss, as a function of the share of deleted edges restored.
  Lines are means of paired gains over ten damage orders. Shading shows one
  standard deviation.}
  \label{fig:repair}
\end{figure}

\begin{figure}[!htbp]
  \centering
  \includegraphics[width=0.95\linewidth]{act3/final/primekg_5000_bfs_v1/fig_restored_edge_types.pdf}
  \caption{Relation types among the restored edges at a 20\% budget. Bars show the
  mean share and error bars one standard deviation over ten damage orders.}
  \label{fig:repair-types}
\end{figure}

These tests capture only the variation between damage orders. All of them use
the same 60 test queries, which behave like about 27 independent ones
(Section~\ref{sec:blind-uncertainty}), and 19 of the 21 test drugs also appear in
training queries. A 4-point difference corresponds to two or three queries.

The best repair therefore depends on the task, in a way the theory explains. A
two-hop query is a two-layer object, so at large budgets the right policy
restores the two layers it uses, and Eq.~\eqref{eq:r2} gives the size of the
gain. When the horizon allows a PPI step, or when the budget is small enough
that only bridges matter, structural importance becomes the better guide.
